\section{Engram Hybrid Adaptive Consensus}
\label{sec:design}

\subsection{Architecture Overview}
\label{subsec:architecture}

Engram is a Cosmos SDK application~\cite{kwon2016cosmos} running on a
forked CometBFT core. Three
external layers are continuously monitored: (i) Bitcoin, used as a
settlement and finality anchor via periodic checkpointing; (ii)
Celestia~\cite{albassam2019lazyledger},
used as the data-availability (DA) layer; and (iii) the validator P2P
network itself. Rather than assuming these layers are always healthy, every
validator runs deterministic local sensors against each of them, and the
resulting health readings are embedded directly into the consensus proposal
(Fig.~\ref{fig:architecture}), so that ``is the periphery healthy'' becomes a
question the validator set answers by quorum, exactly as it does for
transaction validity.

\begin{figure}[h]
\centering
% TODO: thay bang anh that, export tu docs/ARCHITECTURE.md (mermaid) hoac
% ve lai bang TikZ o dang vector cho paper.
\fbox{\parbox{0.9\linewidth}{\centering TODO figure -- Architecture: Engram
execution + Bitcoin settlement + Celestia DA + FSM sensors}}
\caption{System architecture: the Engram validator set, its Bitcoin
settlement anchor, its Celestia DA layer, and the FSM sensors that bridge
them into the consensus proposal.}
\label{fig:architecture}
\end{figure}

\subsection{Threat Model}
\label{subsec:threat-model}

We assume a partially synchronous network: message delay is unknown but
bounded after some unknown global stabilization time, consistent with
CometBFT's own consensus model. The adversary is Byzantine and controls up
to $f$ validators out of $n$, where $f = \lfloor (n-1)/3 \rfloor$ (e.g.,
$f{=}2$ at the $n{=}7$ scale used throughout our evaluation); Byzantine
validators may deviate arbitrarily from the protocol, including colluding
with one another, but cannot forge signatures or break the collision
resistance of the hash functions relied on by the consensus layer, the
Bitcoin blockchain, or the re-anchoring proof circuit, and that the
UltraHonk backend's KZG polynomial-commitment trusted setup was performed
honestly -- standard cryptographic hardness assumptions this paper relies
on but does not independently verify. Beyond the validator set's own Byzantine bound, three
further assumptions bound what the design defends against: (i) the
health-aggregation step that certifies a peripheral-health reading requires
quorum-signed observations, so no single Byzantine validator can force a
false health claim into a proposal on its own; (ii) Bitcoin settlement is
treated as safe against reorganizations no deeper than a configurable
finality depth, consistent with every Bitcoin-anchoring protocol in this
class (Section~\ref{subsec:settlement}); and (iii) honest validators'
peripheral-health views converge within the partial-synchrony bound rather
than being identical at every instant. We additionally assume the FSM's
threshold parameters (e.g.\ the suspicious- and sovereign-state thresholds,
the hysteresis wait, and $K_{\text{DEEP\_FINALITY}}$) are protocol-level
constants that every honest validator runs identically, in the same way
every honest validator runs the same block-time or quorum-size parameter;
this is required for the quorum-agreement mechanism described in
Section~\ref{subsec:fsm} to function at all, since agreement depends on
every honest validator's local evaluation of the transition function
matching. A validator running a divergent
threshold would simply fail to agree with the honest majority and could
not contribute to quorum -- a liveness cost to that validator, not a safety
hazard to the network -- so we do not treat threshold consistency as a
Byzantine attack surface, but we note it here as an operational assumption
this paper does not separately test. A related but distinct scenario is a
genuine split view: validators running identical, correctly-configured
thresholds can still disagree if their own external reads of Bitcoin, DA,
or P2P health differ near a threshold boundary (e.g.\ 3 of 7 validators
observing unhealthy Bitcoin connectivity while 4 observe healthy, neither
side reaching the $\geq 2f{+}1$ quorum needed to commit either target
state). Section~\ref{sec:evaluation} tests decisive splits only (a
self-healing 1-of-7 view and a sub-quorum 2-of-7 view); a persistent split
straddling the quorum boundary has not been tested and could in principle
cause repeated nil-prevotes and round failures rather than a clean commit
either way -- we return to this in Section~\ref{sec:discussion}.
Section~\ref{sec:formal} restates
(i)--(iii) as explicit preconditions on the safety and accountability
property targeted there, and Appendix~\ref{app:reorg} tests assumption
(ii) directly against real Bitcoin reorganizations. We do not assume
Bitcoin, Celestia, or the P2P network are themselves Byzantine-fault
tolerant in the sense the validator set is; the claim this design makes is
that the FSM correctly degrades when any of these external networks
becomes unavailable or unreliable, not that it substitutes for or defends
against a Byzantine minority within Bitcoin's or Celestia's own respective
consensus mechanisms, which is out of scope.

\subsection{Finite State Machine Design}
\label{subsec:fsm}

The protocol maintains a four-state FSM that degrades gracefully across
security levels instead of halting when an external dependency fails
(Table~\ref{tab:fsm-states}). \textsc{Anchored} is normal operation, secured
by both Bitcoin settlement and Celestia DA. \textsc{Suspicious} is an early
warning state, reached after a sustained (not single-block) unhealthy
reading; forced-transaction admission is rate-limited here rather than
withdrawals being touched --
withdrawals remain permitted throughout \textsc{Suspicious}, and are locked
only in \textsc{Sovereign}/\textsc{Recovering} (Table~\ref{tab:fsm-states},
\emph{Circuit Breaker Isolation}, Section~\ref{subsec:formal-invariants}).
\textsc{Sovereign} is an active-partition fallback: the validator set
continues producing blocks under its own quorum, without requiring
external corroboration from Bitcoin or Celestia, and the circuit breaker
locks all cross-chain withdrawals. \textsc{Recovering} is the resolution state, entered
once connectivity is restored, during which all \textsc{Sovereign}-mode
transitions are aggregated into a single recursive zero-knowledge proof for
re-anchoring to Bitcoin.

\begin{table}[h]
\caption{FSM State Comparison}
\label{tab:fsm-states}
\centering
\footnotesize
\begin{tabular}{@{}lllll@{}}
\toprule
\textbf{State} & \textbf{Security basis} & \textbf{Withdrawals} & \textbf{Throughput} & \textbf{Finality} \\
\midrule
\textsc{Anchored} & Bitcoin + Celestia DA & Permitted & Full & $\sim$10--60 min \\
\textsc{Suspicious} & Bitcoin (degraded) & Permitted & Restricted$^{\dagger}$ & $\sim$10--60 min \\
\textsc{Sovereign} & Local quorum & Locked & Full & $\sim$2 s \\
\textsc{Recovering} & Local quorum + proof & Locked & Full & $\sim$2 s \\
\bottomrule
\end{tabular}
\\[2pt]
\footnotesize $^{\dagger}$Forced-transaction admission is rate-limited by a
bounded queue while \textsc{Suspicious} (Section~\ref{subsec:transitions}), not
locked outright. \textsc{Anchored} and \textsc{Suspicious} share the same
underlying Bitcoin finality window -- \textsc{Suspicious} reflects reduced
throughput and admission, not a distinct finality timescale, since Bitcoin
settlement is still degraded rather than replaced. Bitcoin-anchored
finality figures reflect Bitcoin's own confirmation-time distribution and
are not independently measured in this evaluation. More fundamentally,
\textsc{Anchored}'s Bitcoin-derived finality and \textsc{Sovereign}'s BFT
finality are not the same kind of guarantee -- one is probabilistic
economic finality, the other is deterministic BFT finality -- and this
table should not be read as claiming \textsc{Sovereign} is simply
``faster-finalizing'' than \textsc{Anchored}.
\end{table}

Fig.~\ref{fig:fsm} shows the transition diagram. All transitions require
greater than 2/3 quorum agreement through the consensus pipeline: sensors
only \emph{propose} a target state, and a validator prevotes for a proposal
only if the embedded FSM state matches its own local, deterministic
evaluation of the same transition function.

\begin{figure}[h]
\centering
\begin{tikzpicture}[
    >=Stealth, node distance=2.2cm, auto,
    state/.style={rectangle, rounded corners, draw, minimum width=2.4cm,
                  minimum height=1cm, align=center, font=\footnotesize}
]
\node[state] (anchored) {ANCHORED};
\node[state, right=of anchored] (suspicious) {SUSPICIOUS};
\node[state, below=of suspicious] (sovereign) {SOVEREIGN};
\node[state, below=of anchored] (recovering) {RECOVERING};

\draw[->] (anchored) -- node[font=\tiny]{\shortstack{Warning\\($k$ blocks)}} (suspicious);
\draw[->] (anchored) -- node[font=\tiny, swap]{Critical} (sovereign);
\draw[->] (suspicious) -- node[font=\tiny]{\shortstack{Critical /\\timeout}} (sovereign);
\draw[->] (suspicious) to[bend left=15] node[font=\tiny]{\shortstack{Healthy\\($k$ blocks)}} (anchored);
\draw[->] (sovereign) to[bend left=12] node[font=\tiny]{Healthy} (recovering);
\draw[->] (recovering) to[bend left=12] node[font=\tiny]{\shortstack{Critical /\\non-critical$^{\times k}$}} (sovereign);
\draw[->] (recovering) to[bend left=40] node[font=\tiny]{\shortstack{Healthy $\wedge$ safe\_blocks=$H$\\$\wedge$ ZK-proof valid}} (anchored);
\end{tikzpicture}
\vspace{4mm}
\caption{Engram FSM state transition diagram. Every edge into
\textsc{Anchored} or out of \textsc{Suspicious}/\textsc{Recovering} is gated
by a hysteresis counter requiring $k$ (or $H$, the recovery-edge
hysteresis wait) consecutive blocks, not a single reading.}
\label{fig:fsm}
\end{figure}

\subsection{Network Sensors}
\label{subsec:sensors}

Three sensor categories feed the transition function; all are deterministic
so that every honest validator computes the same target state from the same
committed data.

\textbf{Bitcoin finality gap.} $\Delta H_{\text{BTC}} = H_{\text{current}} -
\min(H_{\text{submitted}}, H_{\text{anchored}})$ measures how far the latest
checkpoint is from Bitcoin confirmation. A growing gap indicates Bitcoin
congestion or a liveness attack on checkpointing. Because a delay-only metric
cannot detect a \emph{forged} checkpoint, each validator additionally
performs an off-band SPV check (Merkle-proof inclusion and block-header
verification); a failed SPV check is treated as critical-severity, on par
with an excessive gap, since it means the anchor height is unverifiable
rather than merely stale.

\textbf{Data availability gap.} $\Delta H_{\text{DA}} = H_{\text{local}} -
H_{\text{verified}}$ measures the lag between the current chain head and the
highest block for which a verified DA commitment has been received. The
target design backs this with Celestia-style data-availability sampling and
Blobstream attestation (Celestia's proof-of-inclusion bridge to external
chains); Section~\ref{subsec:sim-to-real} discusses the gap
between that target design and the current prototype, in which each
validator trusts only its own network's Celestia bridge rather than
sampling across the other validators' independent bridges.

\textbf{P2P health (tri-interface profiler).} A raw peer count cannot
distinguish a long-lived honest peer from a freshly injected Sybil, and
eclipse attacks combine several attack vectors that no single metric
captures~\cite{rehman2025eclipse}. Engram's P2P health check instead
combines six structural, behavioral, and temporal bounds -- subnet
diversity, active-anchor-peer count, clean-peer count, peer churn rate,
average peer tenure, and peer latency -- applied uniformly to the Engram
P2P interface, the Bitcoin SPV client, and the Celestia light client.
Section~\ref{subsec:rq1} and Section~\ref{subsec:limitations}
discuss the empirical effectiveness and the acknowledged scope limits of
this profiler, respectively.

\subsection{State Transition Logic and Hysteresis}
\label{subsec:transitions}

A \emph{warning} condition (a suspicious Bitcoin gap, an unhealthy DA
reading, or unhealthy P2P quality) demotes \textsc{Anchored} to
\textsc{Suspicious} only after a sustained streak of consecutive unhealthy
blocks, absorbing single-block noise. A \emph{critical} condition (an
excessive Bitcoin gap, a failed SPV check, complete anchor-peer isolation, or
a suspicious-state timeout) escalates directly to \textsc{Sovereign} from
either \textsc{Anchored} or \textsc{Suspicious}. Recovery in the opposite
direction is deliberately harder than escalation: exiting
\textsc{Suspicious} back to \textsc{Anchored}, or exiting
\textsc{Recovering} back to \textsc{Anchored}, each requires a sustained
streak of healthy blocks under a dedicated hysteresis counter, not a single
healthy reading. This asymmetry -- fast to degrade, slow and quorum-gated to
recover -- is what Section~\ref{subsec:rq3} evaluates
quantitatively across all three hysteresis-gated edges. Separately from
this hysteresis mechanism, \textsc{Suspicious} also caps the rate at which
would-be-censored transactions can be force-included, via a bounded
admission queue rather than by touching withdrawals (which stay permitted
throughout \textsc{Suspicious}, Table~\ref{tab:fsm-states}); this queue cap
is what gives \textsc{Suspicious} its ``Restricted'' throughput entry in
Table~\ref{tab:fsm-states}, and Section~\ref{sec:discussion} returns to it
as a partial mitigation against a repeated-trigger attack on the
\textsc{Anchored}$\leftrightarrow$\textsc{Suspicious} boundary.

\subsection{Extended Proposal and Consensus Protocol}
\label{subsec:proposal}

The base consensus proposal is extended with four additional fields: the
target FSM state, a DA receipt (last DA-verified height and an attestation
flag), a Bitcoin receipt (checkpoint height and canonical block hash), and a
reference to a re-anchoring ZK proof when applicable. A validator casts a
prevote for a proposal only if all of the following hold: the embedded FSM
state matches its own local transition function; the embedded health flag
matches its own aggregated health certificate; the DA receipt and Bitcoin
receipt are internally consistent and within a round-adaptive tolerance
whenever the state is \textsc{Anchored}/\textsc{Recovering} or the
corresponding layer is currently healthy; and, when the proposal claims
exit from \textsc{Recovering} to \textsc{Anchored}, a valid ZK re-anchoring
proof is attached. Section~\ref{subsec:rq4} measures the resulting
proposal-size and validation-cost overhead directly.

\subsection{Re-anchoring via Zero-Knowledge Proof of Recovery}
\label{subsec:reanchoring}

When the chain exits \textsc{Sovereign} mode, re-executing every
locally-produced block on Bitcoin would be prohibitively expensive. Instead,
a single dynamic-padding circuit (Noir, proved with Barretenberg's
UltraHonk backend)~\cite{aztec_noir_barretenberg} attests to header-chain
continuity, FSM-transition legality, and withdrawal-lock adherence across
a \textsc{Sovereign} interval of up to $N_{\max}=256$ headers, at constant
cost regardless of how many of those slots are real versus padding
(Section~\ref{sec:evaluation}). We evaluated, and rejected, a true
recursive-composition alternative in the style of proof-carrying
data~\cite{bitansky2013recursive} -- aggregating many small per-header
proofs via in-circuit verification of each inner proof -- because the
in-circuit cost of verifying one inner UltraHonk proof dwarfs the cost of
directly hashing a header, making recursion 650--1,200$\times$ slower than
the dynamic-padding circuit at matched header counts in our benchmarks
(Appendix~\ref{app:backend}); recursion is a better fit for proofs that
are themselves expensive to re-verify (e.g.\ rollup batches), not cheap
per-header hashing. A \textsc{Sovereign} interval longer than
$N_{\max}$ headers is covered by a rolling sequence of dynamic-padding
proofs, each starting where the last left off, rather than by a single
proof over the whole interval -- so verification is $O(1)$ per proof, not
$O(1)$ regardless of total interval length. Verification itself happens
off-chain: every validator runs an embedded UltraHonk verifier locally, and
only a commitment to the resulting state root is anchored on Bitcoin via
the existing checkpoint mechanism (Section~\ref{subsec:sensors}) -- Bitcoin
itself cannot evaluate the proof (its script is not expressive enough for
UltraHonk verification), so its role after re-anchoring is timestamping
and settlement of that commitment, not independent proof verification. A
full backend comparison against an alternative STARK-based prover
(Plonky3) is deferred to Appendix~\ref{app:backend}, since it evaluates a
secondary design question (prover backend choice) rather than a core
contribution.